% BEGIN REV09_TRANSPORTE_POLAR
\subsection{Transport polaire de l’incidence et conservation des lecteurs}
\label{corpus9:iv:polar}

La réalisation finie du complexe orienté d’incidence
reçoit la frontière et l’orientation construites dans l’état
APP--TRIT--TPK. Une fois fixés les produits scalaires de
cette réalisation, soient \(\mathcal H_0,\mathcal H_1\) ses
espaces finis et
\[
 d:\mathcal H_0\longrightarrow\mathcal H_1,\qquad
 \Delta_0=d^*d,\qquad \Delta_1=dd^*.
\]
Les lecteurs d’incidence et de retour sont
\[
 \rho_{\mathrm{int}}(v)=(dv,d^*dv),\qquad
 \rho_{\mathrm{ref}}(a)=(d^*a,dd^*a).
\]
Les symboles \(\Delta_0,\Delta_1\) désignent ici les
laplaciens de cette réalisation hilbertienne. L’orientation,
les métriques et le domaine de \(d\) font partie de leurs données.

\begin{proposition}[Équivalence unitaire des secteurs positifs
]
\label{corpus9:iv:polar-transporte}
Soient
\[
 \mathcal H_0^+=(\ker d)^\perp,\qquad
 \mathcal H_1^+=(\ker d^*)^\perp=\operatorname{Ran}d.
\]
L'application
\begin{equation}
 \mathcal U=d(\Delta_0|_{\mathcal H_0^+})^{-1/2}
    :\mathcal H_0^+\longrightarrow\mathcal H_1^+
 \label{corpus9:iv:polar-operador}
\end{equation}
est unitaire et satisfait, dans ces secteurs,

\begin{equation}
 \mathcal U\Delta_0=\Delta_1\mathcal U,\qquad
 \mathcal U f(\Delta_0)=f(\Delta_1)\mathcal U
 \label{corpus9:iv:polar-funcional}
\end{equation}
pour toute fonction \(f\) définie sur le spectre positif
fini commun.

\end{proposition}
\begin{proof}
La décomposition en valeurs singulières réduite est \(d=V\Sigma W^*\),
où les colonnes de \(W\) et \(V\) sont des bases orthonormales
de \(\mathcal H_0^+\) et \(\mathcal H_1^+\), et \(\Sigma\)
est diagonale à valeurs strictement positives. Par conséquent,
\[
 \Delta_0|_{\mathcal H_0^+}=W\Sigma^2W^*,\qquad
 \Delta_1|_{\mathcal H_1^+}=V\Sigma^2V^*,\qquad
 \mathcal U=VW^*.
\]
Il s’ensuit \(\mathcal U^*\mathcal U=I_{\mathcal H_0^+}\) et
\(\mathcal U\mathcal U^*=I_{\mathcal H_1^+}\). Évaluer \(f\)
sur la même diagonale \(\Sigma^2\) prouve
\eqref{corpus9:iv:polar-funcional}. Si \(d=0\), les deux
secteurs positifs sont nuls et les identités ont leur
interprétation unique entre espaces de dimension nulle.
\end{proof}

Pour \(v\in\mathcal H_0^+\) et \(h,t\geq0\), on conserve
l’énergie quadratique, la résolvante et le semi-groupe :

\begin{align}
 \langle\mathcal Uv,\Delta_1\mathcal Uv\rangle
    &=\langle v,\Delta_0v\rangle,\\
 \mathcal U(I+h\Delta_0)^{-1}
    &=(I+h\Delta_1)^{-1}\mathcal U,\\
 \mathcal Ue^{-t\Delta_0}
    &=e^{-t\Delta_1}\mathcal U.
 \label{corpus9:iv:polar-lectores}
\end{align}
La première identité se déduit de l’unitarité et de
l’entrelacement. Les deux autres sont des cas de
\eqref{corpus9:iv:polar-funcional}, avec
\(f(x)=(1+hx)^{-1}\) et \(f(x)=e^{-tx}\).


\subsubsection{Réalisation sur le cycle tritique}
Dans les bases orthonormées des sommets et des arêtes orientées
du cycle à trois éléments,
\[
 d=
 \begin{pmatrix}
 -1&1&0\\
 0&-1&1\\
 1&0&-1
 \end{pmatrix}.
\]
La multiplication des matrices donne
\[
 \Delta_0=\Delta_1=
 \begin{pmatrix}
 2&-1&-1\\
 -1&2&-1\\
 -1&-1&2
 \end{pmatrix}
 =3I-J_3=2I-C-C^{-1},
\]
où \(J_3\) a toutes ses entrées égales à un et \(C\)
est la permutation cyclique. Le vecteur constant engendre le
noyau ; sur son complément orthogonal, \(J_3=0\).
Ainsi, le spectre est \(\{0,3,3\}\) et
\(\mathcal U=d/\sqrt3\) dans le secteur positif.
Les deux lecteurs partagent ce spectre parce qu’ils proviennent de la
même incidence, en conservant leurs espaces de sommets et
d’arêtes.

\subsubsection{Orientation et compatibilité des raffinements
}
La substitution \(d\mapsto-d\) conserve \(\Delta_0,\Delta_1\)
et change \(\mathcal U\) en \(-\mathcal U\). Plus
généralement, pour une inversion d’orientations d’arêtes
représentée par une matrice diagonale unitaire \(E\),
\[
 d'=Ed,\qquad \Delta_0'=\Delta_0,\qquad
 \Delta_1'=E\Delta_1E^*,\qquad \mathcal U'=E\mathcal U.
\]
Ces égalités s’obtiennent en utilisant \(E^*E=I\) dans les
définitions. L’orientation est ainsi conservée par le
transporteur, même lorsque les lectures spectrales
coïncident.


Dans une collection de raffinements, les applications
\(J_m^{(0)}:\mathcal H_{0,m}^+\to\mathcal H_{0,m+1}^+\)
et
\(J_m^{(1)}:\mathcal H_{1,m}^+\to\mathcal H_{1,m+1}^+\)
doivent conserver les domaines et les normes et satisfaire
\[
 \mathcal U_{m+1}J_m^{(0)}
    =J_m^{(1)}\mathcal U_m
\]
pour transporter cette identification à travers les niveaux.
La proposition prouve l’équivalence finie des secteurs
positifs ; le carré précédent exprime la condition supplémentaire
du raffinement. Les noyaux, les restrictions non linéaires
et les passages à la limite conservent leurs démonstrations propres.

% END REV09_TRANSPORTE_POLAR
